%% -*- coding: utf-8 -*-
%% Discussion section for gPAC paper

\section{Discussion}

\subsection{Summary of Findings}
We presented TorchPAC, a GPU-accelerated framework for phase-amplitude coupling computation built on PyTorch \cite{paszke_pytorch_2019}. Our benchmarking results demonstrated that GPU parallelization reduced PAC computation time by approximately two orders of magnitude compared to TensorPAC \cite{combrisson_tensorpac_2020}, a widely used CPU-based implementation. This speedup was consistent across variations in batch size, number of channels, recording duration, and sampling frequency, indicating that the performance gain generalizes across typical experimental configurations rather than being limited to narrow parameter regimes. Importantly, the computed PAC values showed close agreement with those from TensorPAC, as quantified by mean squared error, confirming that the acceleration did not compromise computational accuracy. Additionally, by implementing all operations as differentiable PyTorch modules, the framework supports gradient-based optimization of filter parameters, enabling data-driven selection of frequency bands for phase and amplitude extraction.

\subsection{Relation to Existing Tools}
Several software packages exist for PAC computation. TensorPAC \cite{combrisson_tensorpac_2020} provides a flexible Python interface supporting multiple PAC metrics, including the modulation index \cite{tort_measuring_2010} and the mean vector length \cite{canolty_high_2006}, but relies on CPU execution and processes channels sequentially. PACTools and similar packages share this sequential processing limitation. Our framework addresses this bottleneck by exploiting the inherent parallelism in PAC computation---bandpass filtering, Hilbert transformation, and modulation index calculation can be applied independently across channels, trials, and frequency band pairs. While TensorPAC remains suitable for small-scale analyses where computation time is not a constraint, TorchPAC is designed for settings where the number of channels, trials, or frequency band combinations makes sequential processing impractical.

\subsection{Implications for Neuroscience Research}
The reduction in computation time has practical consequences for PAC-based analyses. First, it enables systematic exploration of coupling across many frequency band combinations, which is necessary for studies examining cross-frequency interactions without strong prior hypotheses about which bands are involved. Second, the processing speed approaches levels required for online PAC estimation, which is relevant for brain-computer interface applications where PAC-derived features may serve as control signals \cite{canolty_high_2006}. Third, the trainable parameter functionality allows frequency band boundaries to be optimized jointly with downstream classification objectives, potentially reducing reliance on manually specified frequency ranges that may not capture subject-specific coupling patterns. Our demonstration of 95\% classification accuracy on a synthetic coupled-versus-uncoupled discrimination task provides initial evidence that this end-to-end optimization is feasible, though validation on physiological data with clinical or cognitive labels remains necessary.

\subsection{Limitations}
Several limitations should be noted. First, the framework requires access to a CUDA-compatible GPU; users without such hardware would need to run in CPU mode, forfeiting the speed advantage. Second, GPU memory imposes an upper bound on the size of data that can be processed in a single batch. While our chunked processing strategy mitigates this constraint, very large multi-electrode recordings may still require careful memory management. Third, our speed benchmarking was conducted on a single hardware configuration (NVIDIA GeForce RTX 4090), and the magnitude of speedup will vary with GPU model and CPU baseline. Fourth, the current validation against TensorPAC was limited to the modulation index metric; equivalence for other PAC measures (e.g., mean vector length, phase-locking value) has not yet been established. Finally, the trainable PAC module was evaluated only on synthetic data with known ground-truth coupling, and its utility for identifying physiologically relevant frequency bands in empirical recordings requires further investigation.

\subsection{Future Directions}
Several extensions of this work merit exploration. Evaluating the trainable PAC module on clinical datasets---such as intracranial recordings from patients with epilepsy, where PAC abnormalities have been documented \cite{tort_measuring_2010}---would test whether data-driven frequency band selection improves detection of pathological coupling. Implementing the framework within real-time acquisition pipelines would clarify its viability for closed-loop neuromodulation and brain-computer interface applications. Extending the approach to related cross-frequency coupling measures, including phase-phase coupling and amplitude-amplitude coupling, would broaden its applicability. Finally, systematic benchmarking across a range of GPU architectures would provide users with practical guidance for hardware selection.

\label{sec:discussion}

%%%% EOF
